% ── §9b: Detailed Statistical Framework ──
\section{Detailed Statistical Framework}
\label{sec:statistical-framework}

\subsection{Hypothesis Testing for Contact Enrichment}

For each candidate pair $(i,j)$, we observe a binary outcome: the pair is
either a native contact ($Y = 1$) or not ($Y = 0$).  Under the null
hypothesis that our scoring method ranks pairs randomly, the number of
contacts among the top-$K$ predictions follows a hypergeometric distribution.
When $K \ll N_{\mathrm{cand}}$, this is well-approximated by a binomial:
\begin{equation}
	X \sim \mathrm{Binomial}(K, \pi_0), \qquad
	\pi_0 = \frac{N_c}{N_{\mathrm{cand}}},
\end{equation}
where $\pi_0$ is the base contact rate.  We reject the null at significance
level $\alpha$ when $P(X \geq x_{\mathrm{obs}} \mid K, \pi_0) < \alpha$
(one-sided test for enrichment).

For the Gray-adjacency test (PATH~E), the same binomial framework applies
with $\pi_0$ replaced by the background rate of Gray-adjacent pairs among
non-contacting residue pairs, matched within each protein and separation bin
to control for compositional and phylogenetic confounding.

\subsection{Permutation Tests for Encoding Specificity}

To assess whether the He/Gray encoding specifically (rather than any encoding
of amino acids into 20 slots) produces Gray-adjacency enrichment, we perform
a label-permutation test.  For each of $R = 30$ permutations $\sigma$, we
reassign codewords according to $\sigma: \{0,\ldots,19\} \to \{0,\ldots,19\}$
and recompute all pairwise Gray-Hamming distances under the permuted codes.
The permutation $p$-value is
\begin{equation}
	p_\mathrm{perm} =
	\frac{1 + |\{r : E_r \geq E_\mathrm{obs}\}|}{1 + R},
\end{equation}
where $E_\mathrm{obs}$ is the observed enrichment and $E_r$ is the enrichment
under permutation $r$.  The resulting $z$-score quantifies how many standard
deviations the true encoding exceeds the random-encoding distribution.

\subsection{Bootstrap Confidence Intervals}

Protein-level bootstrap resampling provides confidence intervals that account
for between-protein heterogeneity.  For $B$ bootstrap replicates, we draw
proteins with replacement from the set of $P$ targets and recompute the
statistic of interest on each bootstrap sample.  The 2.5th and 97.5th
percentiles of the bootstrap distribution give a 95\% confidence interval.
This procedure is used for the Gray-adjacency enrichment (PATH~E) with $B =
	2000$.

\subsection{Wilcoxon Signed-Rank Test for Paired Comparisons}

When comparing two methods evaluated on the same set of test proteins, we use
the Wilcoxon signed-rank test on the paired differences.  This non-parametric
test does not assume normality and is appropriate for precision values that
may be skewed toward zero for methods with sparse top-$K$ hits.  The null
hypothesis is that the median difference between paired observations is zero;
the alternative is that it differs (two-sided) or exceeds/falls below zero
(one-sided).

\subsection{Multiple Comparisons}

When testing multiple hypotheses simultaneously, we apply Bonferroni
correction to control the family-wise error rate.  For $m$ tests, the
per-test significance threshold becomes $\alpha / m$.  In practice, the
strongest findings in this work have $p$-values many orders of magnitude below
any corrected threshold, while weaker effects are explicitly flagged as not
surviving correction.
